\chapter{Sieves, Grothendieck Topologies, and Sites}\label{ch03:chap}

\section{Purpose and scope}\label{ch03:sec:purpose}

This chapter develops the covering apparatus: sieves, their
identification with subfunctors of representables, pullback sieves,
Grothendieck topologies and sites, coverage bases (pretopologies) and the
topology they generate, and subcanonicity. As with Chapter~2, everything
here is standard material, presented at tutorial depth and cited to the
standard sources --- Mac Lane--Moerdijk \cite{maclane}, the Stacks
Project \cite{stacks00ZC,stacks00W1,stacks00WO}, and Johnstone
\cite{johnstone} --- and none of it is a project contribution. The
chapter's role in the monograph is to put on the table, with complete
proofs, the small collection of elementary facts on which everything in
Parts~I and~II turns: that sieves on a poset are down-closed sets of
elements; that a sieve containing an identity, or a split epimorphism,
is maximal (the two facts that drive the legacy collapse of Chapter~5);
that a pretopology generates a topology (the passage that legitimizes
the successor's generating families in Chapter~9); and that the union
coverage on a powerset lattice is a subcanonical Grothendieck topology
(the ambient standard fact of which Chapter~9's join theorem identifies
the project-specific instance).

Two complete worked enumerations anchor the abstractions: all sieves on
each object of the four small posets of Chapter~2, and all Grothendieck
topologies on the two-element chain. The second enumeration is
deliberately exhaustive --- there are exactly four --- because
exhaustive enumeration over a small structure is the method of the
finite certificate of Chapter~11, and seeing the method by hand on the
smallest nontrivial site calibrates what the certificate's larger
enumerations do and do not establish.

\section{Sieves}\label{ch03:sec:sieves}

\subsection{Definition and first examples}\label{ch03:sec:sieve-def}

\begin{definition}[Sieve]\label{ch03:def:sieve}
Let $\cC$ be a category and $C \in \Ob(\cC)$. A \emph{sieve} on $C$ is
a set $S$ of morphisms, all with codomain $C$, that is closed under
precomposition: if $f : D \to C$ belongs to $S$ and $h : E \to D$ is
any morphism of $\cC$, then $f \circ h$ belongs to $S$.
\end{definition}

The mental picture: a sieve on $C$ is a family of ``ways of approaching
$C$'' that is closed downward --- any way of approaching $C$ that
factors through an admitted way is itself admitted. The empty set is a
sieve (the closure condition is vacuous), and the set of \emph{all}
morphisms into $C$ is a sieve; the latter is important enough to name.

\begin{definition}[Maximal sieve]\label{ch03:def:maximal}
The \emph{maximal sieve} on $C$, written $m_C$, is the set of all
morphisms of $\cC$ with codomain $C$. It contains $\id_C$, and it is
the largest sieve on $C$.
\end{definition}

\begin{definition}[Generated sieve]\label{ch03:def:generated}
For a family $\mathcal{F}$ of morphisms with codomain $C$, the
\emph{sieve generated by} $\mathcal{F}$, written
$\langle \mathcal{F} \rangle$, is
\[
\langle \mathcal{F} \rangle
= \{\, f \circ h \mid f \in \mathcal{F},\ h \text{ a morphism with }
\mathrm{cod}(h) = \mathrm{dom}(f) \,\}.
\]
\end{definition}

\begin{proposition}\label{ch03:prop:generated}
$\langle \mathcal{F} \rangle$ is the smallest sieve on $C$ containing
$\mathcal{F}$.
\end{proposition}

\begin{proof}
It contains $\mathcal{F}$, since $f = f \circ \id$. It is a sieve: a
precomposition $(f \circ h) \circ h' = f \circ (h \circ h')$ is again
of the required form, by associativity. And any sieve $T$ containing
$\mathcal{F}$ contains every $f \circ h$ with $f \in \mathcal{F}$, by
$T$'s own closure; hence $\langle \mathcal{F} \rangle \subseteq T$.
\end{proof}

\subsection{Sieves on a poset}\label{ch03:sec:sieve-poset}

On a poset the notion collapses to something completely concrete.
Recall that $\mathord{\downarrow} V = \{U : U \le V\}$ denotes the
principal down-set of $V$, and call a subset
$D \subseteq \mathord{\downarrow} V$ \emph{down-closed} if
$U' \le U \in D$ implies $U' \in D$.

\begin{proposition}[Sieves on a poset are down-closed
sets]\label{ch03:prop:sieve-downset}
Let $P$ be a poset and $V \in P$. The assignment
$S \mapsto \{\, U \in P \mid \iota_{U,V} \in S \,\}$ is a bijection
between sieves on $V$ (in the sense of
Definition~\ref{ch03:def:sieve}) and down-closed subsets of
$\mathord{\downarrow} V$. Under this bijection the maximal sieve
corresponds to $\mathord{\downarrow} V$ itself, the empty sieve to
$\varnothing$, and the sieve generated by a family
$\{U_i \le V\}_{i \in I}$ to the union of principal down-sets
$\bigcup_{i \in I} \mathord{\downarrow} U_i$.
\end{proposition}

\begin{proof}
A morphism with codomain $V$ is exactly a $\iota_{U,V}$ with
$U \le V$, and distinct $U$ give distinct morphisms, so a set of
morphisms into $V$ is exactly a subset of
$\mathord{\downarrow} V$. Precomposition closure translates as: if
$U \in S$ (meaning $\iota_{U,V} \in S$) and $U' \le U$, then
$\iota_{U,V} \circ \iota_{U',U} = \iota_{U',V} \in S$, i.e.\
$U' \in S$ --- precisely down-closedness. The three identifications
follow by unwinding: $m_V$ collects all $U \le V$; the generated
sieve of Definition~\ref{ch03:def:generated} collects all $U'$
factoring through some $U_i$, i.e.\
$\bigcup_i \mathord{\downarrow} U_i$.
\end{proof}

From here on, when working on a poset we identify a sieve with its
down-closed set of elements and write $U \in S$ without comment.

\begin{example}[All sieves on the small posets]\label{ch03:ex:sieve-enum}
Using Proposition~\ref{ch03:prop:sieve-downset}, sieves on an object
$V$ are down-closed subsets of $\mathord{\downarrow} V$, and we can
enumerate them completely for the four posets of Chapter~2.
\begin{enumerate}
\item \emph{Two-element chain} $[2]$ ($0 \le 1$). On $0$:
$\mathord{\downarrow} 0 = \{0\}$, whose down-closed subsets are
$\varnothing$ and $\{0\}$ --- two sieves, the empty and the maximal.
On $1$: $\mathord{\downarrow} 1 = \{0, 1\}$; down-closed subsets are
$\varnothing$, $\{0\}$, and $\{0, 1\}$ --- three sieves. (Note
$\{1\}$ alone is \emph{not} down-closed: $0 \le 1$ would force
$0 \in$.) Total over both objects: $5$.
\item \emph{Three-element chain} $[3]$. On $2$: down-closed subsets
of $\{0,1,2\}$ are $\varnothing, \{0\}, \{0,1\}, \{0,1,2\}$ ---
four sieves. On $1$: three. On $0$: two. Total: $9$.
\item \emph{The vee} $\Lambda$ ($x, y \le t$). On $t$:
$\mathord{\downarrow} t = \{x, y, t\}$; down-closed subsets:
$\varnothing$, $\{x\}$, $\{y\}$, $\{x, y\}$, $\{x, y, t\}$ --- five
sieves. (Any set containing $t$ must contain $x$ and $y$.) On $x$:
two. On $y$: two. Total: $9$. These five sieves on $t$ are exactly
the five subpresheaves of $yt$ enumerated in Chapter~2
(Example~\ref{ch02:ex:subrep-vee}) --- the coincidence is
Proposition~\ref{ch03:prop:sieve-subfunctor} below.
\item \emph{The square} $B = \mathcal{P}(\{1,2\})$. On the top
$\{1,2\}$: down-closed subsets of the whole four-element lattice.
Any such set containing the top is everything; the sets omitting the
top are the down-closed subsets of the three-element sub-poset
$\{\varnothing, \{1\}, \{2\}\}$, namely $\varnothing$,
$\{\varnothing\}$, $\{\varnothing, \{1\}\}$,
$\{\varnothing, \{2\}\}$, $\{\varnothing, \{1\}, \{2\}\}$ --- five
of them --- so there are six sieves on the top. On $\{1\}$:
$\mathord{\downarrow}\{1\} = \{\varnothing, \{1\}\}$ is a
two-element chain, giving three sieves; likewise on $\{2\}$. On
$\varnothing$: two. Total: $6 + 3 + 3 + 2 = 14$.
\end{enumerate}
The same computation scales to the sixteen-element context lattice
$L = \mathcal{P}(\Omega)$ of Part~II; the total there is $298$
sieves over the sixteen contexts, a count derived in closed form in
Chapter~7 (Lemma~\ref{ch07:lem-downsets} and
Theorem~\ref{ch07:thm-298}), tabulated in Appendix~B, and
independently corroborated by the exhaustive enumeration in the
registered certificate.
\end{example}

\begin{exercise}\label{ch03:exr:sieve-count-chain}
Show that on an $n$-element chain the number of sieves on the top
element is $n + 1$, and the total number of sieves over all objects
is $2 + 3 + \cdots + (n+1) = \frac{(n+1)(n+2)}{2} - 1$. Verify the
formula against the enumerations above for $n = 2, 3$.
\end{exercise}

\subsection{Sieves as subfunctors of representables}\label{ch03:sec:sieve-subfunctor}

\begin{proposition}[Sieve--subfunctor correspondence]\label{ch03:prop:sieve-subfunctor}
Let $\cC$ be a small category and $C \in \Ob(\cC)$. For a sieve $S$
on $C$, define
\[
\widetilde{S}(D) = \{\, f \in \Hom_{\cC}(D, C) \mid f \in S \,\}
\subseteq yC(D),
\]
with restriction maps inherited from $yC$ (precomposition). Then
$S \mapsto \widetilde{S}$ is a bijection between sieves on $C$ and
subpresheaves of the representable $yC$.
\end{proposition}

\begin{proof}
First, $\widetilde{S}$ is a subpresheaf precisely because $S$ is a
sieve: the restriction of $yC$ along $h : E \to D$ sends
$f \mapsto f \circ h$, and the requirement that this map carry
$\widetilde{S}(D)$ into $\widetilde{S}(E)$
(Definition~\ref{ch02:def:subpresheaf}) says exactly that
$f \in S$ implies $f \circ h \in S$ --- the closure condition of
Definition~\ref{ch03:def:sieve}. So the assignment is well defined,
and conversely, any subpresheaf $T \subseteq yC$ yields the set
$S_T = \bigcup_D T(D)$ of morphisms into $C$, which is
precomposition-closed by the same translation read backwards.

The two assignments are mutually inverse: starting from a sieve $S$,
the set $S_{\widetilde{S}}$ collects the elements of all the
$\widetilde{S}(D)$, which is $S$ itself, since a morphism into $C$
lies in exactly one $\Hom_{\cC}(D, C)$ (namely for $D$ its domain);
starting from a subpresheaf $T$, the subpresheaf
$\widetilde{S_T}$ has $\widetilde{S_T}(D) = S_T \cap
\Hom_{\cC}(D, C) = T(D)$.
\end{proof}

This correspondence is what lets sieves interact with presheaves: a
\emph{matching family} for a sieve, defined below
(Definition~\ref{ch03:def:matching}), is exactly a natural
transformation $\widetilde{S} \Rightarrow F$, and the sheaf condition
of Chapter~4 compares such transformations with sections over $C$ in
Yoneda style. On a poset, the correspondence specializes to the
coincidence already observed: sieves on $V$ = down-closed subsets of
$\mathord{\downarrow} V$ = subpresheaves of the down-set indicator
$yV$ (Proposition~\ref{ch02:prop:representable-poset}).

\subsection{Pullback sieves}\label{ch03:sec:pullback-sieves}

\begin{definition}[Pullback sieve]\label{ch03:def:pullback-sieve}
Let $S$ be a sieve on $C$ and $g : D \to C$ a morphism. The
\emph{pullback sieve} $g^* S$ on $D$ is
\[
g^* S = \{\, h : E \to D \mid g \circ h \in S \,\}.
\]
\end{definition}

$g^* S$ is a sieve: if $g \circ h \in S$ and $k$ is composable, then
$g \circ (h \circ k) = (g \circ h) \circ k \in S$ by closure of $S$.
Intuitively, $g^* S$ is the sieve $S$ ``as seen from $D$ through
$g$'': an approach to $D$ is admitted iff its composite into $C$ was
admitted.

\begin{lemma}[Membership via pullback]\label{ch03:lem:membership}
For any sieve $S$ on $C$ and morphism $g : D \to C$:
\begin{enumerate}
\item $g^* m_C = m_D$;
\item $g \in S$ if and only if $g^* S = m_D$;
\item if $S \subseteq T$ then $g^* S \subseteq g^* T$.
\end{enumerate}
\end{lemma}

\begin{proof}
(1) Every $h$ into $D$ satisfies $g \circ h \in m_C$. (2) If
$g \in S$, then for every $h : E \to D$ we have
$g \circ h \in S$ by closure, so $h \in g^* S$; thus
$g^* S = m_D$. Conversely if $g^* S = m_D$ then
$\id_D \in g^* S$, which says $g \circ \id_D = g \in S$. (3) is
immediate from the definition.
\end{proof}

\begin{example}[Pullback sieves on a poset]\label{ch03:ex:pullback-poset}
On a poset, with $S$ a down-closed subset of
$\mathord{\downarrow} V$ and $U \le V$, the pullback along
$U \to V$ is
\[
(\iota_{U,V})^* S = \{\, W \le U \mid W \in S \,\}
= S \cap \mathord{\downarrow} U :
\]
the part of $S$ lying below $U$. For instance, on
$B = \mathcal{P}(\{1,2\})$, let $S$ be the sieve on the top
generated by $\{1\}$, i.e.\
$S = \mathord{\downarrow} \{1\} = \{\varnothing, \{1\}\}$. Its
pullback along $\{2\} \subseteq \{1,2\}$ is
$S \cap \mathord{\downarrow} \{2\} = \{\varnothing\}$: from the
vantage of $\{2\}$, only the empty overlap of the admitted approach
survives. Its pullback along $\{1\} \subseteq \{1,2\}$ is
$\{\varnothing, \{1\}\} = m_{\{1\}}$, the maximal sieve ---
consistent with Lemma~\ref{ch03:lem:membership}(2), since
$\{1\} \in S$.
\end{example}

\section{Grothendieck topologies}\label{ch03:sec:topologies}

\subsection{The axioms}\label{ch03:sec:axioms}

\begin{definition}[Grothendieck topology; site]\label{ch03:def:topology}
A \emph{Grothendieck topology} $J$ on a category $\cC$ assigns to
each object $C$ a collection $J(C)$ of sieves on $C$, called
\emph{covering sieves}, such that:
\begin{enumerate}
\item[(J1)] (\emph{Maximality}) $m_C \in J(C)$ for every $C$;
\item[(J2)] (\emph{Stability}) if $S \in J(C)$ and $g : D \to C$ is
any morphism, then $g^* S \in J(D)$;
\item[(J3)] (\emph{Transitivity, or local character}) if
$S \in J(C)$ and $T$ is any sieve on $C$ such that
$f^* T \in J(D)$ for every $f : D \to C$ in $S$, then
$T \in J(C)$.
\end{enumerate}
A \emph{site} is a pair $(\cC, J)$ of a category and a topology on
it. See \cite[Ch.~III]{maclane} and \cite{stacks00Z4}.
\end{definition}

Each axiom encodes one piece of covering common sense, and we take
them one at a time, because the successor construction of
Part~II is engineered directly against them.

\emph{J1, maximality.} The family of all ways into $C$ covers $C$.
This is the trivial admission: an object is always covered by
everything mapping into it, including its own identity. J1 is
harmless --- but it is also the axiom through which the legacy
collapse operates: since maximal sieves are \emph{automatically}
covering in every topology, a clause whose generated sieves are all
maximal contributes literally nothing to a topology. Being maximal
is the covering-theoretic way of being redundant.

\emph{J2, stability.} Covers restrict. If $S$ covers $C$ and we
change vantage to $D$ along $g : D \to C$, the restricted family
$g^* S$ must cover $D$. In bookkeeping terms: admission of a family
of local views is not allowed to depend on staying at the original
target; any sub-target inherits an admitted family. On a poset,
stability says: if the down-closed set $S$ covers $V$, then
$S \cap \mathord{\downarrow} U$ covers every $U \le V$. This is the
axiom that forces the successor's generating families to be
pullback-stable, which is the content of the coverage-basis lemma
of Chapter~9 and is where the intersection formula
$(V \cap B) \cap U = U \cap B$ does its work.

\emph{J3, transitivity.} Covers compose hierarchically. If $S$
covers $C$, and $T$ is a sieve that covers \emph{locally on} $S$
--- pulled back along each member of $S$, it covers that member's
domain --- then $T$ covers $C$ outright. This is the axiom that
lets one refine covers: cover $C$ by pieces, cover each piece by
smaller pieces, and conclude that the assembly of smaller pieces
covers $C$. In Part~II it is precisely J3 that makes the
\emph{common refinement} of the ACBE and RPM covers --- the
four-singleton family --- covering for the joined topology, which
is the join theorem of Chapter~9.

\subsection{First consequences}\label{ch03:sec:consequences}

\begin{lemma}[Monotonicity]\label{ch03:lem:monotone}
Let $J$ be a Grothendieck topology, $S \in J(C)$, and $T$ a sieve on
$C$ with $S \subseteq T$. Then $T \in J(C)$.
\end{lemma}

\begin{proof}
We verify the hypothesis of (J3) for $T$ against the covering sieve
$S$. Let $f : D \to C$ be a member of $S$. Then $f \in T$ as well,
so by Lemma~\ref{ch03:lem:membership}(2), $f^* T = m_D$, which is
covering by (J1). Thus $f^* T \in J(D)$ for every $f \in S$, and
(J3) yields $T \in J(C)$.
\end{proof}

\begin{lemma}[Intersection stability]\label{ch03:lem:intersection}
If $S, T \in J(C)$ then $S \cap T \in J(C)$.
\end{lemma}

\begin{proof}
$S \cap T$ is a sieve (both closure conditions hold). First
observe that pullback commutes with intersection:
$h \in f^*(S \cap T)$ iff $f \circ h \in S$ and
$f \circ h \in T$, i.e.\ $f^*(S \cap T) = f^* S \cap f^* T$. Now
let $f : D \to C$ be any member of $S$. Then
$f \in S$ gives $f^* S = m_D$
(Lemma~\ref{ch03:lem:membership}(2)), so
$f^* (S \cap T) = f^* T$, which is covering by (J2) applied to
$T$. Thus the hypothesis of (J3) holds for $S \cap T$ against $S$,
and $S \cap T \in J(C)$.
\end{proof}

Together, Lemmas~\ref{ch03:lem:monotone}
and~\ref{ch03:lem:intersection} say each $J(C)$ is a \emph{filter}
of sieves: upward closed and closed under finite intersection ---
a useful sanity check when enumerating topologies, as we do in
Section~\ref{ch03:sec:enumeration}.

\subsection{The two decisive elementary facts}\label{ch03:sec:decisive}

The next two lemmas are, mathematically, the least sophisticated
results in this monograph; they are also the ones with the largest
downstream consequences, since Chapter~5's collapse theorem is
nothing but their application to the legacy clauses. We therefore
prove them with full care.

\begin{lemma}[A sieve containing an identity is
maximal]\label{ch03:lem:id-maximal}
Let $S$ be a sieve on $C$ with $\id_C \in S$. Then $S = m_C$.
\end{lemma}

\begin{proof}
Let $f : D \to C$ be any morphism with codomain $C$. Since
$\id_C \in S$ and $S$ is closed under precomposition
(Definition~\ref{ch03:def:sieve}), the composite
$\id_C \circ f$ belongs to $S$. But $\id_C \circ f = f$ by
unitality (C2). Hence every morphism into $C$ belongs to $S$, i.e.\
$S = m_C$.
\end{proof}

\begin{lemma}[A sieve containing a split epimorphism is
maximal]\label{ch03:lem:splitepi-maximal}
Let $S$ be a sieve on $C$ containing a split epimorphism
$\varepsilon : B \to C$, say with section $\eta : C \to B$,
$\varepsilon \circ \eta = \id_C$. Then $S = m_C$.
\end{lemma}

\begin{proof}
By precomposition closure applied to $\varepsilon \in S$ and the
morphism $\eta : C \to B$, the composite
$\varepsilon \circ \eta$ belongs to $S$. By hypothesis
$\varepsilon \circ \eta = \id_C$, so $\id_C \in S$, and
Lemma~\ref{ch03:lem:id-maximal} gives $S = m_C$.
\end{proof}

\begin{remark}\label{ch03:rem:decisive-consequence}
Combined with the observation under (J1) above --- maximal sieves
are covering in \emph{every} topology, so they never enlarge one
--- these lemmas give a general design constraint on sites: a
covering clause whose generating families contain an identity, or
a split epimorphism onto the target, is vacuous. It adds no sieve
that was not already covering, and, since topology generation is a
closure operator, it leaves no trace of having been imposed.
Chapter~5 applies this constraint verbatim: the legacy ACBE clause
contains an identity (the reflexive descent map at equal contexts)
and the legacy RPM clause contains a split epimorphism (the counit
$\varepsilon_X$, split by $\eta_X$), so both are vacuous, and the
legacy topology collapses to its observational clause. Note also
Remark~\ref{ch02:rem:no-split-in-posets}: in a poset the only
split epimorphisms are identities, so on the successor's poset
site the second mechanism cannot even be formulated nontrivially
--- a structural immunity, not a lucky accident.
\end{remark}

\begin{exercise}\label{ch03:exr:iso-maximal}
Show more generally that a sieve containing any isomorphism is
maximal. (Every isomorphism is a split epimorphism.)
\end{exercise}

\section{Coverage bases and generated topologies}\label{ch03:sec:bases}

Topologies are rarely specified sieve by sieve; one specifies a
\emph{basis} of covering families and closes under the axioms. The
successor construction of Part~II is specified exactly this way ---
two-element partition families at each context --- so we develop
the passage carefully.

\begin{definition}[Pretopology / coverage basis]\label{ch03:def:pretopology}
Let $\cC$ be a category with pullbacks. A \emph{pretopology} (or
\emph{coverage basis}) $K$ assigns to each object $C$ a collection
$K(C)$ of families of morphisms $\{f_i : C_i \to C\}_{i \in I}$,
called \emph{generating covers}, such that:
\begin{enumerate}
\item[(B1)] (\emph{Isomorphisms}) every family consisting of a
single isomorphism $\{f : C' \xrightarrow{\ \cong\ } C\}$ ---
in particular the identity family $\{\id_C\}$ --- is in $K(C)$;
\item[(B2)] (\emph{Stability}) if
$\{f_i : C_i \to C\}_i \in K(C)$ and $g : D \to C$ is any
morphism, then the pulled-back family
$\{\, \pi_i : C_i \times_C D \to D \,\}_i$ is in $K(D)$;
\item[(B3)] (\emph{Composition}) if
$\{f_i : C_i \to C\}_i \in K(C)$ and for each $i$ a family
$\{g_{ij} : C_{ij} \to C_i\}_j \in K(C_i)$ is given, then the
composite family $\{f_i \circ g_{ij}\}_{i,j}$ is in $K(C)$.
\end{enumerate}
See \cite{stacks00ZC} and \cite[Ch.~III]{maclane} for variants;
weaker ``coverage'' axioms
suffice to generate a topology, but the pretopology form is the one
the successor's families satisfy and is simplest to state.
\end{definition}

\begin{theorem}[Generated topology]\label{ch03:thm:generated-topology}
Let $K$ be a pretopology on a category $\cC$ with pullbacks. Define,
for each object $C$,
\[
J_K(C) = \{\, S \text{ a sieve on } C \mid S \supseteq R
\text{ for some } R \in K(C) \,\},
\]
where $S \supseteq R$ means $S$ contains every member of the family
$R$. Then $J_K$ is a Grothendieck topology --- the \emph{topology
generated by} $K$ --- and it is the smallest topology in which every
generating cover generates a covering sieve.
\end{theorem}

\begin{proof}[Proof sketch, with the key steps in full]
\emph{(J1).} The identity family $\{\id_C\}$ lies in $K(C)$ by
(B1), and $m_C$ contains it; so $m_C \in J_K(C)$.

\emph{(J2).} Let $S \in J_K(C)$ via $R = \{f_i : C_i \to C\}
\subseteq S$, and let $g : D \to C$. By (B2) the pulled-back family
$R' = \{\pi_i : C_i \times_C D \to D\}$ is in $K(D)$. We claim
$R' \subseteq g^* S$: each $\pi_i$ satisfies
$g \circ \pi_i = f_i \circ \pi'_i$ (the pullback square), and
$f_i \circ \pi'_i \in S$ because $f_i \in S$ and $S$ is a sieve;
hence $g \circ \pi_i \in S$, i.e.\ $\pi_i \in g^* S$. So
$g^* S \supseteq R' \in K(D)$, giving $g^* S \in J_K(D)$.

\emph{(J3).} Let $S \in J_K(C)$ via $R = \{f_i\} \subseteq S$, and
let $T$ be a sieve with $f^* T \in J_K(\mathrm{dom} f)$ for every
$f \in S$. In particular, for each $i$, $f_i^* T \supseteq R_i$
for some $R_i = \{g_{ij} : C_{ij} \to C_i\} \in K(C_i)$. Form the
composite family $\{f_i \circ g_{ij}\}_{i,j}$, which lies in
$K(C)$ by (B3). Each of its members lies in $T$: indeed
$g_{ij} \in f_i^* T$ means precisely $f_i \circ g_{ij} \in T$.
Hence $T \supseteq \{f_i \circ g_{ij}\}_{i,j} \in K(C)$, so
$T \in J_K(C)$.

\emph{Minimality.} If $J'$ is any topology in which each generating
cover $R \in K(C)$ generates a covering sieve
$\langle R \rangle \in J'(C)$, then every $S \supseteq R$ satisfies
$S \supseteq \langle R \rangle$ (Proposition~\ref{ch03:prop:generated})
and hence $S \in J'(C)$ by monotonicity
(Lemma~\ref{ch03:lem:monotone}, which used only the axioms). So
$J_K \subseteq J'$ objectwise. Conversely each
$\langle R \rangle$ itself contains $R$, so
$\langle R \rangle \in J_K(C)$: the generated topology does make
every basis cover covering. Full details in the general (coverage)
setting are in \cite{stacks00ZC}; the poset case, where pullbacks
are intersections and all of the above becomes finite bookkeeping,
is rechecked concretely in Chapter~9 for the successor's families.
\end{proof}

\begin{remark}[Sheaf condition on a basis]\label{ch03:rem:basis-sheaf}
The companion fact, stated here and used from Chapter~4 onward: a
presheaf is a sheaf for the generated topology $J_K$ if and only if
it satisfies the gluing condition for the basis families themselves
\cite{stacks00ZC,maclane}. This is what allows Chapter~9 to verify
subcanonicity and sheafhood by checking only the two-piece
partition families rather than every covering sieve of $J_A$ and of
$J_R$ separately (the figure $168$ quoted in
Remark~\ref{ch03:rem:ch9-preview} counts the covering sieves of the
\emph{join} $J_A \vee J_R$; each factor is strictly coarser and has
strictly fewer); the
certificate of Chapter~11, by contrast, computes the full
topologies and rechecks the equivalence exhaustively --- a
deliberate redundancy between basis-level proof and
topology-level computation.
\end{remark}

\section{Matching families and subcanonicity}\label{ch03:sec:subcanonical}

The full theory of sheaves belongs to Chapter~4, but subcanonicity
is a property \emph{of topologies} and is most naturally introduced
here. We give the minimal forward definitions.

\begin{definition}[Matching family; gluing]\label{ch03:def:matching}
Let $F$ be a presheaf on $\cC$ and $S$ a sieve on $C$. A
\emph{matching family} for $S$ with values in $F$ is an assignment,
to each $f : D \to C$ in $S$, of an element $x_f \in F(D)$, such
that
\[
F(h)(x_f) = x_{f \circ h}
\quad \text{for every } h : E \to D .
\]
Equivalently (Proposition~\ref{ch03:prop:sieve-subfunctor}), it is
a natural transformation $\widetilde{S} \Rightarrow F$. Every
section $s \in F(C)$ induces a matching family by restriction:
$x_f = F(f)(s)$. The presheaf $F$ is a \emph{sheaf} for a topology
$J$ if for every $C$ and every $S \in J(C)$ this canonical map
from $F(C)$ to matching families is a bijection; it is
\emph{separated} if the map is always injective.
\end{definition}

\begin{definition}[Subcanonical topology]\label{ch03:def:subcanonical}
A topology $J$ on $\cC$ is \emph{subcanonical} if every
representable presheaf $yC$ is a $J$-sheaf. The \emph{canonical
topology} is the largest subcanonical topology; every topology
contained in a subcanonical one is subcanonical
\cite{stacks00WO,johnstone}.
\end{definition}

Subcanonicity is the license to regard site objects as evidence
objects: it says the topology's compatibility demands never
manufacture or destroy maps into a fixed target, so $C$ and $yC$
can be identified inside the sheaf category. When it fails, the
identification is unearned --- the exact situation of the legacy
site, where subcanonicity is not derivable from the registered
axioms (Chapter~6), and the reason the program's representation
rule sheafifies legacy objects before use. On the successor's
context factor, by contrast, subcanonicity is a theorem
(Chapter~9), and Section~\ref{ch03:sec:union} below proves the
standard prototype of that theorem in its natural generality.

\begin{example}[A sheaf check on the square, by hand]\label{ch03:ex:sheaf-preview}
On $B$, let $S = \{\varnothing, \{1\}, \{2\}\}$ be the sieve on the
top generated by the two singletons, and let $G$ be the bit
presheaf (Example~\ref{ch02:ex:bit-presheaf}). A matching family
assigns $x_{\{1\}} \in G(\{1\})$, $x_{\{2\}} \in G(\{2\})$,
$x_\varnothing \in G(\varnothing)$ with all restrictions
compatible; since $G(\varnothing)$ is a singleton, the only
constraints are $x_{\{1\}}|_\varnothing = x_\varnothing =
x_{\{2\}}|_\varnothing$, which hold automatically. So matching
families correspond to free pairs
$(x_{\{1\}}, x_{\{2\}})$ --- there are $2 \times 2 = 4$ --- and
the canonical map from $G(\{1,2\})$, which has $4$ elements, sends
a bit-function on $\{1,2\}$ to its two restrictions. This map is a
bijection (a function on $\{1,2\}$ is exactly a pair of values), so
$G$ satisfies the sheaf condition for this sieve. The
constant-verdict presheaf $H$
(Example~\ref{ch02:ex:verdict-presheaf}) does not: matching
families again form a $2 \times 2$ set, but $H(\{1,2\})$ has only
$2$ elements, and the canonical map hits only the diagonal. The
count $4$ versus $2$ is the germ of the witness argument of
Chapter~9, where the same arithmetic ($2^2 = 4$ sections against a
singleton matching-family set, and conversely) separates the ACBE
and RPM topologies.
\end{example}

\section{All Grothendieck topologies on the two-element chain}\label{ch03:sec:enumeration}

We now carry out, in full, the smallest interesting classification:
every Grothendieck topology on the poset $[2]$
($0 \le 1$). The answer is that there are exactly four, and the
proof method --- identify the free boolean choices, derive forcing
implications among them, then verify each consistent assignment ---
is the hand-scale version of the fixed-point saturation the
certificate of Chapter~11 performs on the sixteen-object lattice.

Recall the sieve inventory (Example~\ref{ch03:ex:sieve-enum}): on
$0$, the empty sieve $\varnothing_0$ and the maximal
$m_0 = \{0\}$; on $1$, the empty sieve $\varnothing_1$, the
proper sieve $S_f = \{0\}$ (the down-set generated by $0$,
corresponding to the morphism $f : 0 \to 1$), and the maximal
$m_1 = \{0, 1\}$.

By (J1), any topology contains $m_0$ and $m_1$. Three boolean
choices remain:
\[
a = [\varnothing_0 \in J(0)], \qquad
b = [S_f \in J(1)], \qquad
c = [\varnothing_1 \in J(1)].
\]

\begin{lemma}[Forcing implications]\label{ch03:lem:forcing}
In any Grothendieck topology $J$ on $[2]$:
\begin{enumerate}
\item if $\varnothing_C \in J(C)$ for an object $C$, then
\emph{every} sieve on $C$ is in $J(C)$;
\item $c \Rightarrow a$;
\item $c \Rightarrow b$;
\item $(a \wedge b) \Rightarrow c$.
\end{enumerate}
\end{lemma}

\begin{proof}
(1) Apply (J3) with $S = \varnothing_C$: the hypothesis ``for
every $f \in S$, $f^* T$ is covering'' quantifies over the empty
set and is vacuously true for every sieve $T$; hence every $T$ is
covering. (2) If $\varnothing_1 \in J(1)$, then by (J2) its
pullback along $f : 0 \to 1$ is covering; and
$f^* \varnothing_1 = \{h \mid f \circ h \in \varnothing_1\} =
\varnothing_0$. So $a$ holds. (3) Immediate from (1) applied at
$C = 1$: if $c$, then all sieves on $1$ are covering, in
particular $S_f$. (4) Assume $a$ and $b$. Apply (J3) with
$S = S_f \in J(1)$ and $T = \varnothing_1$. The members of
$S_f$, as morphisms, are exactly $f$ (the unique morphism
$0 \to 1$; its precompositions with $\id_0$ give $f$ again). For
this member, $f^* \varnothing_1 = \varnothing_0$, which is
covering by $a$. So the hypothesis of (J3) holds, and
$\varnothing_1 = T \in J(1)$, i.e.\ $c$.
\end{proof}

\begin{theorem}[Classification]\label{ch03:thm:four-topologies}
The poset $[2]$ carries exactly four Grothendieck topologies:
\begin{enumerate}
\item $J_{\min}$: $(a, b, c) = (0, 0, 0)$ ---
$J(0) = \{m_0\}$, $J(1) = \{m_1\}$;
\item $J_{\mathrm{dense}}$: $(0, 1, 0)$ ---
$J(0) = \{m_0\}$, $J(1) = \{m_1, S_f\}$;
\item $J_{\mathrm{closed}}$: $(1, 0, 0)$ ---
$J(0) = \{m_0, \varnothing_0\}$, $J(1) = \{m_1\}$;
\item $J_{\max}$: $(1, 1, 1)$ --- every sieve on every object is
covering.
\end{enumerate}
\end{theorem}

\begin{proof}
\emph{Only these four are possible.} The eight boolean assignments
$(a, b, c)$ are cut down by Lemma~\ref{ch03:lem:forcing}:
$(0,0,1)$, $(0,1,1)$ violate (2); $(1,0,1)$ violates (3);
$(1,1,0)$ violates (4). The remaining four are those listed.

\emph{Each of the four is a topology.} (J1) holds in all four by
construction. We check (J2) and (J3) case by case; throughout,
pullback along an identity is the identity operation, and pullback
of a maximal sieve is maximal (Lemma~\ref{ch03:lem:membership}(1)),
so only nonmaximal sieves and the morphism $f$ need attention.

$J_{\min}$: (J2): only maximal sieves are covering, and their
pullbacks are maximal. (J3): let $S = m_1$, and $T$ with
$g^* T \in J$ for all $g \in m_1$; taking $g = \id_1$ gives
$T = \id_1^* T \in J(1)$, so the conclusion holds trivially; the
same argument at $0$, and with $S = m_0$. Valid.

$J_{\mathrm{dense}}$: (J2): $S_f$ pulled back along $\id_1$ is
$S_f \in J(1)$; along $f : 0 \to 1$ it is
$f^* S_f = \{h \mid f \circ h \in S_f\} \ni \id_0$ (since
$f \in S_f$), hence $= m_0 \in J(0)$
(Lemma~\ref{ch03:lem:membership}(2)). (J3): suppose
$S \in J(1)$ and $T$ has all pullbacks covering. If
$\id_1 \in S$ (i.e.\ $S = m_1$), then taking $g = \id_1$ gives
$T \in J(1)$ directly. If $S = S_f$, the members are $f$ and its
precompositions (all equal to $f$), and the hypothesis gives
$f^* T \in J(0) = \{m_0\}$, so $\id_0 \in f^* T$, i.e.\
$f \in T$. A sieve on $1$ containing $f$ is either $S_f$ or
$m_1$, and both are in $J(1)$; so $T$ is covering. At the object
$0$ the argument of $J_{\min}$ applies. Valid.

$J_{\mathrm{closed}}$: (J2): the only nonmaximal covering sieve is
$\varnothing_0$ on $0$; morphisms into $0$ are $\id_0$ only, and
$\id_0^* \varnothing_0 = \varnothing_0 \in J(0)$. (J3) at $0$:
covered by Lemma~\ref{ch03:lem:forcing}(1)'s reasoning ---
$\varnothing_0$ covering makes every sieve on $0$ required to be
covering, and indeed both sieves on $0$ \emph{are} covering in
$J_{\mathrm{closed}}$, so no contradiction; for $S = m_0$ the
identity argument applies. (J3) at $1$: the only covering sieve is
$m_1$; for $S = m_1$ and any $T$ with all pullbacks covering,
$g = \id_1$ gives $T \in J(1)$ as before --- note the hypothesis
for $T = S_f$ would require $\id_1^* S_f = S_f \in J(1)$, which
fails, so nothing forces $S_f$ or $\varnothing_1$ in. Valid.

$J_{\max}$: (J2) and (J3) hold trivially since every sieve is
covering. Valid.
\end{proof}

\begin{remark}[Reading the four]\label{ch03:rem:four-reading}
The four topologies illustrate the design space in miniature.
$J_{\min}$ admits only trivial covers: presheaves and sheaves
coincide. $J_{\mathrm{dense}}$ admits the proper sieve $S_f$: it
declares that ``$0$ covers $1$,'' and its sheaves are exactly the
presheaves whose restriction $F(1) \to F(0)$ is a bijection
(matching families for $S_f$ are single elements of $F(0)$, since
the compatibility conditions are trivial; the gluing map is the
restriction). $J_{\mathrm{closed}}$ admits the empty sieve on
$0$: its sheaves must have $F(0)$ a singleton (the empty sieve has
exactly one matching family --- the empty one --- so the sheaf
condition forces $|F(0)| = 1$), with $F(1)$ arbitrary.
$J_{\max}$ forces both conditions: only the terminal presheaf
survives. The correspondence of these four with the four
subtoposes of the presheaf topos on $[2]$ is an instance of the
standard bijection between Grothendieck topologies and
Lawvere--Tierney topologies \cite[Ch.~V]{maclane}
\cite{johnstone}; we cite it for orientation and do not use it.
For the monograph's purposes the important specimen is
$J_{\mathrm{dense}}$: a \emph{proper} generating morphism
producing a genuinely nontrivial topology whose sheaf condition
has real content --- the two-object preview of what the
successor's proper partition pieces do on $L$, and a first
concrete instance of a nontrivial modality (the sheafification
for $J_{\mathrm{dense}}$ replaces $F(1)$ by $F(0)$).
\end{remark}

\begin{exercise}\label{ch03:exr:discrete-two}
Classify all Grothendieck topologies on the \emph{discrete}
two-object poset (no relation between the objects). (Each object
independently either admits only its maximal sieve or admits both
its sieves; there are no cross-object constraints, so there are
four topologies here as well --- but for a structurally opposite
reason. Compare Example~\ref{ch02:ex:psh-discrete}.)
\end{exercise}

\begin{exercise}\label{ch03:exr:three-chain}
Carry out the same program on the three-element chain $[3]$:
enumerate the boolean choices (one per nonmaximal sieve per
object), derive the forcing implications from (J2) and (J3), and
count the topologies. The bookkeeping is heavier; the method is
identical.
\end{exercise}

\section{The union coverage on a powerset lattice}\label{ch03:sec:union}

The final section develops the standard example that Part~II's
context factor instantiates: the union coverage on a powerset
lattice $\mathcal{P}(X)$, for a finite set $X$. This is the
``geometric'' covering notion --- a family of subsets covers $V$
when it actually covers $V$, elementwise --- and it is the
topology that the successor's joined topology $J_A \vee J_R$ turns
out to equal on $L = \mathcal{P}(\Omega)$ (the join theorem of
Chapter~9, ledger \textsc{thm-nt2-join-discrete-context}). Proving
its properties here, in the general finite-powerset setting and
with standard citations, keeps Chapter~9 free to focus on what is
project-specific (the two coarser partition topologies and their
join) and keeps the boundary between textbook fact and project
candidate sharp.

\begin{definition}[Union coverage]\label{ch03:def:union-coverage}
On the poset $\mathcal{P}(X)$, let $K_{\cup}(V)$ consist of all
families $\{V_i \subseteq V\}_{i \in I}$ of subsets of $V$ (each
$V_i \to V$ the inclusion) such that $\bigcup_{i \in I} V_i = V$.
Note the empty family qualifies at $V = \varnothing$ (an empty
union is empty), so the empty context is covered by the
\emph{null family}.
\end{definition}

\begin{proposition}\label{ch03:prop:union-pretopology}
$K_{\cup}$ is a pretopology on $\mathcal{P}(X)$.
\end{proposition}

\begin{proof}
Pullbacks exist and are intersections
(Corollary~\ref{ch02:cor:pullback-powerset}). (B1): the singleton
family $\{V \subseteq V\}$ has union $V$. (B2): pulling
$\{V_i\}_i$ back along $U \subseteq V$ gives the family
$\{V_i \cap U\}_i$, and
$\bigcup_i (V_i \cap U) = \bigl(\bigcup_i V_i\bigr) \cap U
= V \cap U = U$
by distributivity of intersection over union; so the pulled-back
family is in $K_{\cup}(U)$. (B3): if $\bigcup_i V_i = V$ and for
each $i$, $\bigcup_j W_{ij} = V_i$, then
$\bigcup_{i,j} W_{ij} = \bigcup_i V_i = V$. All three axioms hold.
\end{proof}

\begin{definition}\label{ch03:def:union-topology}
The \emph{union topology} $J_{\cup}$ on $\mathcal{P}(X)$ is the
topology generated by $K_{\cup}$
(Theorem~\ref{ch03:thm:generated-topology}).
\end{definition}

\begin{proposition}[Characterization of covering
sieves]\label{ch03:prop:union-characterization}
Let $S$ be a sieve on $V \in \mathcal{P}(X)$, identified with a
down-closed set of subsets of $V$. The following are equivalent:
\begin{enumerate}
\item $S \in J_{\cup}(V)$;
\item $\bigcup_{W \in S} W = V$;
\item $\{x\} \in S$ for every $x \in V$.
\end{enumerate}
In particular, a sieve on $\varnothing$ is covering iff it is
either of the two sieves on $\varnothing$ --- both the empty sieve
and the maximal sieve satisfy (2) vacuously --- so the empty
context is covered even by the empty sieve.
\end{proposition}

\begin{proof}
(1) $\Rightarrow$ (2): if $S$ contains a family
$\{V_i\}$ with $\bigcup_i V_i = V$, then the union of \emph{all}
members of $S$ contains $\bigcup_i V_i = V$, and is contained in
$V$ since every member is a subset of $V$.

(2) $\Rightarrow$ (3): let $x \in V$. By (2), $x \in W$ for some
$W \in S$. Then $\{x\} \subseteq W$, and $S$ is down-closed, so
$\{x\} \in S$.

(3) $\Rightarrow$ (1): the family of singletons
$\{\{x\} \subseteq V\}_{x \in V}$ has union $V$, hence lies in
$K_{\cup}(V)$; by (3) it is contained in $S$, so $S$ is covering
by definition of the generated topology. (For $V = \varnothing$
this is the null family, and every sieve on $\varnothing$
contains it trivially.)
\end{proof}

Condition (3) exhibits $J_{\cup}$ as generated by the
\emph{singleton partition} at each context --- the finite
powerset lattice's atomic decomposition. This is exactly the form
in which the joined topology appears in Chapter~9: the common
refinement of the ACBE and RPM partitions is the singleton
partition (transversality), and the join theorem identifies
$J_A \vee J_R$ with the $J_{\cup}$ of the sixteen-element
lattice. The seed artifacts' internal label $J_{\mathrm{atom}}$
for that join is a name for this topology; as Chapter~1 noted,
the publication layer prefers $J_{\mathrm{singleton}}$, and
neither label asserts the standard category-theoretic ``atomic
topology,'' which is a different notion.

\begin{theorem}[The union topology is
subcanonical]\label{ch03:thm:union-subcanonical}
Every representable presheaf on $\mathcal{P}(X)$ is a sheaf for
$J_{\cup}$. Hence $J_{\cup}$ is subcanonical, as is every
topology contained in it.
\end{theorem}

\begin{proof}
Let $yW$ be a representable
(Proposition~\ref{ch02:prop:representable-poset}: $yW(V)$ is a
singleton if $V \subseteq W$, else empty). By
Remark~\ref{ch03:rem:basis-sheaf} it suffices to check the gluing
condition for basis families. Let
$\{V_i \subseteq V\}_i \in K_{\cup}(V)$, so
$\bigcup_i V_i = V$.

\emph{Case $V \subseteq W$.} There is exactly one section over
$V$, and its restrictions form the unique candidate matching
family; each $V_i \subseteq V \subseteq W$, so each
$yW(V_i)$ is a singleton and there is exactly one matching family.
Bijection: one section, one family.

\emph{Case $V \not\subseteq W$.} There are no sections over $V$,
so we must show there is no matching family either. Suppose one
existed: then $yW(V_i) \neq \varnothing$ for every $i$, i.e.\
$V_i \subseteq W$ for every $i$; but then
$V = \bigcup_i V_i \subseteq W$, contradiction. So both sides are
empty, and the gluing map is (vacuously) a bijection.

\emph{The null family at $\varnothing$.} A matching family over
the empty family is the empty assignment (exactly one), and
$yW(\varnothing)$ is a singleton since
$\varnothing \subseteq W$ always. Bijection again.

The final clause holds because a sheaf for a finer topology is a
sheaf for any coarser one --- the coarser topology's covering
sieves are a subset of the finer's, so the sheaf condition is a
weaker demand. See \cite{stacks00WO} for the general context of
representable sheaves and \cite{johnstone} for the canonical
topology on a frame, of which this is the finite powerset
instance.
\end{proof}

\begin{remark}[What Chapter~9 adds]\label{ch03:rem:ch9-preview}
Everything in this section is standard. What is project-specific,
and deferred to Chapter~9, is the \emph{pair of coarser
topologies}: the ACBE topology $J_A$ and RPM topology $J_R$
generated by the two transverse two-block partitions of
$\Omega$, each strictly coarser than $J_{\cup}$, each
subcanonical (by the final clause above, once contained in
$J_{\cup}$ --- though Chapter~9 proves subcanonicity directly
from the basis as well), each supporting a nontrivial modality,
with join exactly $J_{\cup}$. The union topology is the ambient
standard object; the successor's contribution, whatever its
eventual assessment, lies in the choice and analysis of the
coarser pair, and the certificate's counts ($298$ sieves over
the sixteen contexts, of which $168$ cover for the join,
computed identically from both the join and singleton
presentations) quantify the sixteen-object instance
exhaustively.
\end{remark}

\begin{exercise}\label{ch03:exr:union-square}
On $B = \mathcal{P}(\{1,2\})$, enumerate $J_{\cup}$ completely:
for each of the four objects, list the covering sieves among the
sieve inventory of Example~\ref{ch03:ex:sieve-enum}. (Answers:
on $\{1,2\}$, the covering sieves are
$\{\varnothing, \{1\}, \{2\}\}$ and the maximal sieve --- the
two sieves containing both singletons; on $\{1\}$: the maximal
sieve and $\{\varnothing, \{1\}\}$ coincide, and the covering
sieves are exactly those containing $\{1\}$, i.e.\ the maximal
one; on $\{2\}$: likewise; on $\varnothing$: both sieves. Check
each against
Proposition~\ref{ch03:prop:union-characterization}(3), and count:
$2 + 1 + 1 + 2 = 6$ covering sieves out of $14$.)
\end{exercise}

\section{Summary}\label{ch03:sec:summary}

The chapter's toolkit, as Part~II will draw on it: sieves as
down-closed sets on a poset
(Proposition~\ref{ch03:prop:sieve-downset}) and as subfunctors of
representables (Proposition~\ref{ch03:prop:sieve-subfunctor});
pullback sieves and the membership criterion
(Lemma~\ref{ch03:lem:membership}); the topology axioms with their
covering common sense, and the filter properties
(Lemmas~\ref{ch03:lem:monotone}--\ref{ch03:lem:intersection});
the two decisive maximality facts
(Lemmas~\ref{ch03:lem:id-maximal}
and~\ref{ch03:lem:splitepi-maximal}), which Chapter~5 aims at the
legacy clauses; pretopologies and the generated topology
(Theorem~\ref{ch03:thm:generated-topology}), which Chapter~9 aims
at the successor's partition families; matching families and
subcanonicity (Definitions~\ref{ch03:def:matching}
and~\ref{ch03:def:subcanonical}); the complete classification of
topologies on $[2]$
(Theorem~\ref{ch03:thm:four-topologies}), a hand-scale rehearsal
of exhaustive site analysis; and the union coverage on a finite
powerset lattice with its covering-sieve characterization and
subcanonicity
(Propositions~\ref{ch03:prop:union-pretopology}--\ref{ch03:prop:union-characterization},
Theorem~\ref{ch03:thm:union-subcanonical}), the standard fact of
which the successor's join theorem is the marked, two-coordinate
instance. The whole toolkit is textbook material
\cite{maclane,stacks00ZC,stacks00W1,stacks00WO,johnstone};
Chapter~4 completes the standard background with sheaves and
sheafification, and Chapter~5 begins the project-specific story.
